\documentclass[11pt,a4paper]{article}

\usepackage[a4paper,margin=1.8cm]{geometry}
\usepackage{fontspec}
\usepackage{polyglossia}
\setmainlanguage{russian}
\setotherlanguage{english}
\IfFontExistsTF{DejaVu Serif}{\setmainfont{DejaVu Serif}}{\setmainfont{Times New Roman}}
\IfFontExistsTF{DejaVu Sans}{\setsansfont{DejaVu Sans}}{\setsansfont{Arial}}
\renewcommand{\familydefault}{\sfdefault}
\usepackage{xcolor}
\usepackage{enumitem}
\usepackage{hyperref}
\usepackage{titlesec}
\usepackage{parskip}
\usepackage{multicol}
\usepackage{graphicx}

\hypersetup{
  colorlinks=true,
  urlcolor=[HTML]{2563EB},
  linkcolor=[HTML]{2563EB}
}

\definecolor{accent}{HTML}{2563EB}
\definecolor{textgray}{HTML}{4B5563}
\definecolor{dategray}{HTML}{9CA3AF}
\definecolor{lightgray}{HTML}{E5E7EB}

\setlist[itemize]{leftmargin=1.2em,itemsep=0.25em,topsep=0.35em}
\setlist[enumerate]{leftmargin=1.4em,itemsep=0.25em,topsep=0.35em}
\setlength{\parindent}{0pt}
\setlength{\parskip}{0.45em}

\titleformat{\section}
  {\Large\bfseries\color{accent}}
  {}
  {0pt}
  {}
  [\vspace{0.2em}\titlerule]

\titleformat{\subsection}
  {\large\bfseries}
  {}
  {0pt}
  {}

\newcommand{\roleline}[2]{
  \textbf{#1}\hfill\textcolor{dategray}{\normalfont #2}
}

\newcommand{\project}[5]{
  \subsection*{#1}
  \textit{#2}

  #3

  \textbf{Стек:} #4

  #5
}

\begin{document}

{\fontsize{24}{28}\selectfont\textbf{Наталья Агапова}}\\[0.2em]
{\large Портфолио · Frontend Developer / Web Designer}

\vspace{0.4em}

\href{mailto:agapnatalya004@mail.ru}{agapnatalya004@mail.ru}
\quad|\quad
\href{https://t.me/nhefy}{@nhefy}
\quad|\quad
\href{https://github.com/natagapova}{github.com/natagapova}

\vspace{0.8em}

\textcolor{textgray}{
Frontend-разработчик и дизайнер с 3+ годами опыта на стыке визуала, интерфейсов и AI-инструментов. Собираю лендинги, промо-сайты и продуктовые интерфейсы с фокусом на адаптивность, визуальную подачу и аккуратную реализацию. Умею работать от структуры и референсов до живого интерфейса: продумывать экран, собирать его в Tilda или коде и доводить до понятного результата.
}

\section*{Ключевые навыки}

\begin{multicols}{2}
\begin{itemize}
  \item HTML, CSS, JavaScript, TypeScript
  \item React
  \item Tilda, Zero Block
  \item Figma, FigJam
  \item Adobe Illustrator, Photoshop
  \item Адаптивная верстка, mobile-first
  \item UI для MVP и промо-сайтов
  \item Перевод дизайна из Figma в интерфейс
  \item Формы, интерактивные блоки, базовые анимации
  \item GitHub, GitLab
  \item AI-инструменты в продуктовых процессах
  \item Английский C1+, русский носитель
\end{itemize}
\end{multicols}

\section*{О себе}

Мне интересно собирать цифровые продукты, в которых важны не только код и структура, но и ощущение от взаимодействия. Сильнее всего чувствую себя в задачах, где нужно совместить визуальный вкус, быструю сборку и аккуратный фронтенд. Работаю как в командных проектах, так и самостоятельно: от концепции и экранов до верстки, адаптива и публикации.

\section*{Опыт}

\roleline{Croissan Studio — frontend, design, AI}{апрель 2025 --- настоящее время}

Работаю на стыке дизайна, фронтенда и AI-продакшена. Участвую в клиентских и внутренних проектах студии: от визуальной концепции и экранов до сборки интерфейсов, брендовых материалов и цифровых креативов.

\begin{itemize}
  \item Участвую в UX/UI и визуальной переработке проектов студии.
  \item Работаю с фронтенд-реализацией интерфейсных блоков и промо-страниц.
  \item Подключаюсь к AI-креативу: генерация изображений, видео и аудио для студийных задач.
  \item Сочетаю дизайнерские и технические роли в зависимости от формата проекта.
\end{itemize}

\roleline{AzimovLab — Frontend Team Lead}{февраль 2024 --- сентябрь 2024}

Стартап, разрабатывающий AI-сервис генерации тестов. Руководила фронтенд-командой и участвовала в сборке MVP-интерфейса.

\begin{itemize}
  \item Планировала работу фронтенд-команды и ставила задачи.
  \item Проводила code review и координировала разработку.
  \item Собирала пользовательский интерфейс для MVP.
  \item Интегрировала фронтенд с API и продуктовой логикой.
\end{itemize}

\roleline{Фриланс — frontend / web design}{осень 2022 --- настоящее время}

\begin{itemize}
  \item Верстка лендингов и промо-сайтов на заказ.
  \item Работа с Tilda, статической версткой и визуальными концепциями.
  \item Проекты в сферах культуры, персональных сайтов и образовательных продуктов.
\end{itemize}

\section*{Избранные проекты}

\project
{Croissan Studio}
{\href{https://croissanstudio.ru}{croissanstudio.ru}}
{
Участвую во фронтенд-разработке студийного сайта для AI- и цифровых проектов. Работаю в команде над интерфейсными блоками, адаптивностью и общей визуальной целостностью страницы. Сайт собран как современная студийная витрина с акцентом на кейсы, экспертизу и понятную презентацию услуг.
}
{фронтенд-реализация, адаптивная верстка, работа с контентными блоками, командная разработка}
{
\begin{itemize}
  \item Командная работа над публичным сайтом студии.
  \item Реализация визуально насыщенных секций и кейсовой подачи.
  \item Внимание к типографике, структуре и пользовательскому восприятию страницы.
\end{itemize}
}

\project
{Nina's Space}
{\href{https://ninas-space.vercel.app/}{Сайт} \quad|\quad \href{https://github.com/natagapova/ninas-space}{GitHub}}
{
Некоммерческий персональный сайт, который я сделала по собственной инициативе для близкой подруги — как камерный цифровой проект и творческую практику вне клиентских задач. Внутри — медитации, аффирмации, мини-игры, дыхательный виджет и мультиязычность. Полностью самостоятельная разработка: от идеи и структуры до реализации интерфейса и интерактивных сценариев.
}
{HTML, CSS, JavaScript}
{
\begin{itemize}
  \item Самостоятельная реализация проекта с нуля.
  \item Интерактивность без фреймворков и без сборщика.
  \item Мультиязычность и дополнительные пользовательские сценарии.
  \item Работа с атмосферой, микроопытом и эмоциональным дизайном.
\end{itemize}
}

\project
{Cruel Games}
{\href{https://cruelgames-website.vercel.app/}{Сайт} \quad|\quad \href{https://github.com/natagapova/cruelgames-website}{GitHub}}
{
Промо-сайт театральной постановки, разработанный как статичный, визуально цельный лендинг с акцентом на атмосферу, информацию о спектакле и удобную мобильную подачу. Фриланс-проект, выполненный как отдельная фронтенд-разработка.
}
{HTML, CSS, JavaScript}
{
\begin{itemize}
  \item Полноэкранный hero-блок и промо-подача проекта.
  \item Адаптивная одностраничная верстка.
  \item Кастомная галерея с переключением превью и свайпом на мобильных.
  \item Оптимизация изображений и аккуратная работа со статикой.
\end{itemize}
}

\project
{AI-смена · Иннополис}
{\href{https://aismena.ru}{aismena.ru}}
{
Лендинг летней смены по искусственному интеллекту для школьников 7--11 классов: одностраничник для родителей с фокусом на конверсию в заявки, доверие и полную информацию о программе без звонка менеджеру. Собрала структуру, визуальную подачу и весь сайт на Tilda в сжатые сроки --- от hero и расписания дня до блоков безопасности, стоимости, FAQ и формы записи.
}
{Tilda, Zero Block, адаптивная сборка, формы, подключение домена}
{
\begin{itemize}
  \item Лендинг под заявки: hero, программа, проекты учеников, условия проживания, команда, цена, FAQ.
  \item Таймлайн расписания дня и дифференциация по уровням подготовки.
  \item Адаптивная вёрстка и типографика под все форматы экранов в Tilda.
  \item Форма заявки, юридический футер и публикация на домене aismena.ru.
\end{itemize}
}

\subsection*{Барилова \& Царёв}
\textit{Kingstep Dance Studio · клиентский проект · запуск в 2026}

Сайт для танцевальной студии Kingstep и хореографов Алины Бариловой и Александра Царёва: главная, каталог курсов и страницы занятий с акцентом на визуальную подачу, атмосферу бренда и понятную навигацию. Отвечала за структуру, визуальную концепцию и ключевые экраны; сейчас проект на финальной стадии сборки в Tilda.

\textbf{Стек:} Figma, Tilda, Adobe Illustrator

\begin{itemize}
  \item Собрала логику сайта: главная, раздел курсов и карточки занятий.
  \item Разработала визуальную концепцию в тёмной эстетике студии.
  \item Подготовила макеты, промо-экраны и графику для сборки в Tilda.
  \item Согласовала подачу с брендом Kingstep Dance Studio.
\end{itemize}

\vspace{0.4em}
\begin{center}
  \includegraphics[width=0.48\textwidth,height=0.42\textheight,keepaspectratio]{images/page1.png}%
  \hfill
  \includegraphics[width=0.48\textwidth,height=0.42\textheight,keepaspectratio]{images/page2.png}\\[0.5em]
  \includegraphics[width=0.48\textwidth,height=0.42\textheight,keepaspectratio]{images/page3.png}%
  \hfill
  \includegraphics[width=0.48\textwidth,height=0.42\textheight,keepaspectratio]{images/page4.png}
\end{center}
\begin{center}
  \textcolor{textgray}{\small Ключевые экраны сайта}
\end{center}

\section*{Как я работаю}

\begin{enumerate}
  \item Бриф и понимание задачи: аудитория, цель страницы, ограничения по срокам и формату.
  \item Структура: собираю логику блоков, приоритеты контента и CTA.
  \item Дизайн и сборка: перевожу решение в экран, верстку или Tilda.
  \item Адаптив: отдельно проверяю мобильную версию и ритм страницы.
  \item Полировка: типографика, визуальная целостность, интерактивные элементы, финальные правки.
\end{enumerate}

\section*{Дополнительный опыт}

\begin{itemize}
  \item Продуктовый и бренд-дизайн: интерфейсы, веб-экраны, мобильные экраны, айдентика, базовые бренд-гайды.
  \item AI-инструменты в продакшене: генерация изображений, видео и аудио для клиентских и студийных задач.
  \item Технический бэкграунд: Python, Java, C, C++, TensorFlow, PyTorch, работа с генеративными и CV-моделями.
\end{itemize}

\section*{Образование}

Университет Иннополис, бакалавриат «Прикладной искусственный интеллект», 2026.

\section*{Публикации}

Холодов Я.~А.; Саллум Х.; Агапова Н.~А. «Подходы квантового отжига к взлому шифрования RSA». \textit{Вопросы кибербезопасности}, 2025, №~4, с.~127--133. DOI: \href{https://doi.org/10.21681/2311-3456-2025-4-127-133}{10.21681/2311-3456-2025-4-127-133}.

\vspace{0.4em}

\textcolor{textgray}{\small
(EN) Kholodov, Y.; Salloum, H.; Agapova, N. ``Quantum Annealing Approaches to Breaking RSA Encryption.'' \textit{Voprosy kiberbezopasnosti}, 2025, vol.~4, pp.~127--133.
}

\vfill

\begin{center}
\textcolor{textgray}{Последнее обновление: май 2026}
\end{center}

\end{document}
